\documentclass[aspectratio=169,10pt]{beamer}
% ═══════════════════════════════════════════════════════════════════════════════
% SHARED PREAMBLE — Foundation Models Course (HIT)
% To be \input by each lecture file AFTER \documentclass
% ═══════════════════════════════════════════════════════════════════════════════

% ─────────────────────────────────────────────────────────────────────────────
% BEAMER THEME & BASIC SETUP
% ─────────────────────────────────────────────────────────────────────────────
\usetheme{Madrid}
\usecolortheme{default}

% Remove navigation symbols
\beamertemplatenavigationsymbolsempty

% ─────────────────────────────────────────────────────────────────────────────
% COLORS — HIT Branding
% ─────────────────────────────────────────────────────────────────────────────
\definecolor{hitblue}{RGB}{0,51,102}        % Primary: HIT Navy
\definecolor{hitgold}{RGB}{192,148,0}       % Accent: HIT Gold
\definecolor{hitlight}{RGB}{230,240,250}    % Light background

% Override Beamer palette
\setbeamercolor{primary}{fg=white,bg=hitblue}
\setbeamercolor{secondary}{fg=hitblue,bg=hitlight}
\setbeamercolor{structure}{fg=hitblue}
\setbeamercolor{alerted text}{fg=hitgold}

% ─────────────────────────────────────────────────────────────────────────────
% PACKAGES
% ─────────────────────────────────────────────────────────────────────────────
\usepackage{amsmath}
\usepackage{amssymb}
\usepackage{mathtools}
\usepackage{booktabs}
\usepackage{graphicx}
\usepackage{hyperref}
\usepackage{tikz}
\usepackage{pgfplots}
\usepackage{tcolorbox}
\usepackage{listings}
\usepackage{xcolor}
\usepackage{algorithm}
\usepackage{algpseudocode}

% ─────────────────────────────────────────────────────────────────────────────
% TikZ LIBRARIES
% ─────────────────────────────────────────────────────────────────────────────
\usetikzlibrary{shapes,arrows,positioning,calc,chains,scopes}
\pgfplotsset{compat=1.17}

% ─────────────────────────────────────────────────────────────────────────────
% CODE LISTINGS
% ─────────────────────────────────────────────────────────────────────────────
\lstset{
  basicstyle=\ttfamily\small,
  breaklines=true,
  commentstyle=\color{gray},
  keywordstyle=\color{hitblue},
  stringstyle=\color{hitgold},
  showstringspaces=false,
  backgroundcolor=\color{hitlight},
  frame=single,
  rulecolor=\color{hitblue}
}

% ─────────────────────────────────────────────────────────────────────────────
% TCOLORBOX STYLES
% ─────────────────────────────────────────────────────────────────────────────
\tcbset{
  hitbox/.style={
    colback=hitlight,
    colframe=hitblue,
    fonttitle=\bfseries,
    boxrule=1pt
  },
  alertbox/.style={
    colback=yellow!10,
    colframe=hitgold,
    fonttitle=\bfseries,
    boxrule=1.5pt
  }
}

% ─────────────────────────────────────────────────────────────────────────────
% CUSTOM COMMANDS
% ─────────────────────────────────────────────────────────────────────────────

% Placeholder for figures
\newcommand{\placeholder}[3]{
  \begin{center}
    \fbox{\begin{minipage}[c][#2][c]{#1}
      \centering\small\color{gray}[Figure: #3]
    \end{minipage}}
  \end{center}
}

% Section divider frame
\newcommand{\sectionframe}[2]{
  \begin{frame}[plain]
    \begin{tikzpicture}[remember picture,overlay]
      \fill[hitblue] (current page.south west) rectangle (current page.north east);
      \node[anchor=center, white, font=\Large\bfseries] at (current page.center) {
        \begin{tabular}{c}
          #1 \\[0.3cm]
          {\small\color{hitgold}#2}
        \end{tabular}
      };
    \end{tikzpicture}
  \end{frame}
}

% ─────────────────────────────────────────────────────────────────────────────
% LOAD CUSTOM MACROS
% ─────────────────────────────────────────────────────────────────────────────
% ═══════════════════════════════════════════════════════════════════════════════
% SHARED MACROS — Mathematical Notation for Foundation Models Course
% ═══════════════════════════════════════════════════════════════════════════════

% ─────────────────────────────────────────────────────────────────────────────
% ACTIVATION & LAYER FUNCTIONS
% ─────────────────────────────────────────────────────────────────────────────
\newcommand{\softmax}{\mathrm{softmax}}
\newcommand{\sigmoid}{\mathrm{sigmoid}}
\newcommand{\relu}{\mathrm{ReLU}}
\newcommand{\gelu}{\mathrm{GELU}}
\newcommand{\LayerNorm}{\mathrm{LayerNorm}}
\newcommand{\FFN}{\mathrm{FFN}}
\newcommand{\MHA}{\mathrm{MHA}}
\newcommand{\Attn}{\mathrm{Attention}}

% ─────────────────────────────────────────────────────────────────────────────
% ATTENTION COMPONENTS
% ─────────────────────────────────────────────────────────────────────────────
\newcommand{\query}{Q}
\newcommand{\key}{K}
\newcommand{\val}{V}
\newcommand{\dmodel}{d_{\text{model}}}
\newcommand{\dk}{d_k}
\newcommand{\nheads}{h}

% Scaled dot-product attention formula
\newcommand{\SDPA}{
  \Attn(\query, \key, \val) = \softmax\left(\frac{\query \key^T}{\sqrt{\dk}}\right) \val
}

\newcommand{\CrossAttn}{\text{CrossAttention}}

% ─────────────────────────────────────────────────────────────────────────────
% PROBABILITY & EXPECTATION
% ─────────────────────────────────────────────────────────────────────────────
\newcommand{\E}{\mathbb{E}}
\newcommand{\Prob}{\mathbb{P}}
\newcommand{\KL}{\mathrm{KL}}
\newcommand{\ELBO}{\mathrm{ELBO}}
\newcommand{\given}{\mid}

% ─────────────────────────────────────────────────────────────────────────────
% NORMS & OPERATIONS
% ─────────────────────────────────────────────────────────────────────────────
\newcommand{\norm}[1]{\left\lVert #1 \right\rVert}
\newcommand{\abs}[1]{\left| #1 \right|}
\newcommand{\inner}[2]{\langle #1, #2 \rangle}
\newcommand{\T}{^{\top}}

% ─────────────────────────────────────────────────────────────────────────────
% VECTORS & MATRICES
% ─────────────────────────────────────────────────────────────────────────────
\newcommand{\bvec}[1]{\mathbf{#1}}
\newcommand{\mat}[1]{\mathbf{#1}}
\newcommand{\iid}{\stackrel{\text{i.i.d.}}{\sim}}
\newcommand{\defeq}{\coloneq}

% ─────────────────────────────────────────────────────────────────────────────
% LANGUAGE MODELING
% ─────────────────────────────────────────────────────────────────────────────
\newcommand{\vocab}{\mathcal{V}}
\newcommand{\context}{\text{context}}
\newcommand{\token}{t}
\newcommand{\seq}{x}
\newcommand{\ptheta}{p_\theta}
\newcommand{\pref}{p_{\text{ref}}}

% ─────────────────────────────────────────────────────────────────────────────
% RLHF & DPO
% ─────────────────────────────────────────────────────────────────────────────
\newcommand{\piref}{\pi_{\text{ref}}}
\newcommand{\pitheta}{\pi_\theta}
\newcommand{\yw}{y_w}
\newcommand{\yl}{y_l}
\newcommand{\DPOloss}{\mathcal{L}_{\mathrm{DPO}}}

% ─────────────────────────────────────────────────────────────────────────────
% RETRIEVAL & RAG
% ─────────────────────────────────────────────────────────────────────────────
\newcommand{\docs}{\mathcal{D}}
\newcommand{\qvec}{q}
\newcommand{\dvec}{d}

% ─────────────────────────────────────────────────────────────────────────────
% AGENTS & REASONING
% ─────────────────────────────────────────────────────────────────────────────
\newcommand{\obs}{o}
\newcommand{\act}{a}
\newcommand{\thought}{t}
\newcommand{\tools}{\mathcal{T}}
\newcommand{\history}{h}
\newcommand{\concat}{\oplus}

% ─────────────────────────────────────────────────────────────────────────────
% SETS & LOGIC
% ─────────────────────────────────────────────────────────────────────────────
\newcommand{\R}{\mathbb{R}}
\newcommand{\N}{\mathbb{N}}
\newcommand{\Z}{\mathbb{Z}}

\endinput


% ─────────────────────────────────────────────────────────────────────────────
% TITLE & AUTHOR (can be overridden in individual lectures)
% ─────────────────────────────────────────────────────────────────────────────
\author[D. Lederman]{Dror Lederman}
\institute{Holon Institute of Technology (HIT) \\ Data Science Department}
\date{\today}

% ─────────────────────────────────────────────────────────────────────────────
% FOOTER & HEADER
% ─────────────────────────────────────────────────────────────────────────────
\setbeamertemplate{footline}{
  \leavevmode%
  \hbox{%
    \begin{beamercolorbox}[wd=0.33\paperwidth,ht=2.25ex,dp=1ex,center]{author in head/foot}%
      \usebeamerfont{author in head/foot}\insertshortauthor
    \end{beamercolorbox}%
    \begin{beamercolorbox}[wd=0.34\paperwidth,ht=2.25ex,dp=1ex,center]{title in head/foot}%
      \usebeamerfont{title in head/foot}\insertshorttitle
    \end{beamercolorbox}%
    \begin{beamercolorbox}[wd=0.33\paperwidth,ht=2.25ex,dp=1ex,right]{frame number}%
      \usebeamerfont{frame number}\insertframenumber{} / \inserttotalframenumber\hspace*{2ex}
    \end{beamercolorbox}%
  }%
  \vskip0pt%
}

\endinput


\title{Week 1: From Deep Learning to Foundation Models}
\subtitle{Foundation Models and Agentic AI}
\author{Lecturer Name}
\institute{Holon Institute of Technology (HIT)}
\date{Semester, 2025}

\begin{document}

% ── Title slide ───────────────────────────────────────────────────────────────
\begin{frame}
  \titlepage
\end{frame}

% ── Learning objectives ───────────────────────────────────────────────────────
\begin{frame}{Learning Objectives}
  After this lecture you will be able to:
  \begin{itemize}
    \item Explain the concept of \textbf{transfer learning} and why it transformed NLP.
    \item Describe neural \textbf{scaling laws} and what they predict about model performance.
    \item Define and give examples of \textbf{emergent capabilities} in large models.
    \item Articulate the \textbf{foundation model paradigm} and how it differs from task-specific modeling.
    \item Identify key \textbf{opportunities and risks} associated with foundation models.
  \end{itemize}
\end{frame}

% ══════════════════════════════════════════════════════════════════════════════
\section{From Task-Specific Models to Transfer Learning}
% ══════════════════════════════════════════════════════════════════════════════

\sectionframe{From Task-Specific to Transfer Learning}{How the field moved from bespoke models to shared representations}

\begin{frame}{The Pre-Transfer Learning Era}
  \begin{columns}[T]
    \column{0.48\textwidth}
      \textbf{Classic ML pipeline}
      \begin{itemize}
        \item Hand-craft features (TF-IDF, n-grams)
        \item Train a task-specific classifier
        \item Repeat for every new task
      \end{itemize}
      \vspace{0.5em}
      \textbf{Problems:}
      \begin{itemize}
        \item No knowledge sharing across tasks
        \item Requires large labeled datasets per task
        \item Feature engineering is brittle
      \end{itemize}
    \column{0.48\textwidth}
      \placeholder{0.9\textwidth}{3.5cm}{Classic NLP pipeline: features → classifier per task}
  \end{columns}
  \note{Remind students of their NLP course. Ask: how did we handle sentiment analysis vs NER before 2018?}
\end{frame}

\begin{frame}{Transfer Learning in Computer Vision}
  \begin{itemize}
    \item \textbf{ImageNet pre-training} (2012): AlexNet learned general visual features.
    \item Fine-tuning: freeze early layers, retrain head on target task.
    \item Demonstrated that \emph{features transfer} across tasks.
  \end{itemize}
  \vspace{0.5em}
  \begin{tcolorbox}[hitbox, title=Key Insight]
    Layers in deep networks learn increasingly \textbf{abstract, task-agnostic representations}.
    Early layers $\approx$ task-agnostic; later layers $\approx$ task-specific.
  \end{tcolorbox}
  \vspace{0.3em}
  \textbf{Question:} Can the same idea work for language?
  \note{This motivates BERT / GPT. The answer was yes, but it required rethinking the pretraining objective.}
\end{frame}

\begin{frame}{Transfer Learning Comes to NLP}
  \begin{columns}[T]
    \column{0.55\textwidth}
      \textbf{Key milestones:}
      \begin{enumerate}
        \item \textbf{Word2Vec / GloVe} (2013–2014): static word embeddings
        \item \textbf{ELMo} (2018): contextualized embeddings from LSTMs
        \item \textbf{BERT} (2019): deep bidirectional transformer, fine-tune on downstream tasks~\cite{devlin2019bert}
        \item \textbf{GPT} (2018–2019): left-to-right language model, zero/few-shot~\cite{radford2019gpt2}
      \end{enumerate}
    \column{0.40\textwidth}
      \placeholder{0.95\textwidth}{4cm}{Transfer learning timeline in NLP}
  \end{columns}
  \note{Emphasize the paradigm shift: the same pretrained model, minimal changes → many tasks.}
\end{frame}

% ══════════════════════════════════════════════════════════════════════════════
\section{Scaling Laws}
% ══════════════════════════════════════════════════════════════════════════════

\sectionframe{Neural Scaling Laws}{What happens when you make models bigger?}

\begin{frame}{Scaling Laws: Empirical Findings}
  \textbf{Kaplan et al.\ (2020)}~\cite{kaplan2020scaling} showed that language model loss follows
  \emph{power laws} in three dimensions:
  \vspace{0.4em}
  \begin{columns}[T]
    \column{0.52\textwidth}
      \begin{enumerate}
        \item \textbf{Parameters} $N$: $L(N) \propto N^{-\alpha_N}$
        \item \textbf{Dataset size} $D$: $L(D) \propto D^{-\alpha_D}$
        \item \textbf{Compute} $C$: $L(C) \propto C^{-\alpha_C}$
      \end{enumerate}
      \vspace{0.5em}
      Exponents $\alpha \approx 0.05$–$0.076$ (empirically measured).\\
      \vspace{0.3em}
      \alert{Critical result:} Optimal compute allocation is \emph{not} all-in on parameters.
    \column{0.44\textwidth}
      % tikz/scaling_laws.tex
% Conceptual log-log scaling law curves (compute, parameters, data)
% Usage: % tikz/scaling_laws.tex
% Conceptual log-log scaling law curves (compute, parameters, data)
% Usage: % tikz/scaling_laws.tex
% Conceptual log-log scaling law curves (compute, parameters, data)
% Usage: \input{../../tikz/scaling_laws}
\begin{tikzpicture}
\begin{axis}[
  width=0.82\textwidth,
  height=5.5cm,
  xmode=log, ymode=log,
  xlabel={Scale (compute / parameters / data)},
  ylabel={Loss},
  xlabel style={font=\small},
  ylabel style={font=\small},
  xmin=1e2, xmax=1e10,
  ymin=0.5, ymax=10,
  grid=major,
  grid style={gray!25, very thin},
  legend pos=north east,
  legend style={font=\small},
  tick label style={font=\tiny},
  title style={font=\small\bfseries},
]
% Compute scaling curve
\addplot[hitblue, very thick, domain=1e2:1e10, samples=50]
  {8 * x^(-0.05)};
\addlegendentry{Compute}

% Parameters scaling curve
\addplot[hitgold, very thick, dashed, domain=1e2:1e10, samples=50]
  {7 * x^(-0.045)};
\addlegendentry{Parameters}

% Data scaling curve
\addplot[hitalert, very thick, dotted, domain=1e2:1e10, samples=50]
  {9 * x^(-0.06)};
\addlegendentry{Dataset size}

% Marker: GPT-2 region
\node[font=\tiny, hitblue, rotate=30] at (axis cs:1e4, 3.8) {GPT-2};
% Marker: GPT-3 region
\node[font=\tiny, hitblue, rotate=30] at (axis cs:1e8, 1.8) {GPT-3};
% Marker: Chinchilla-optimal
\draw[->, gray, thin] (axis cs:3e7,2.5) -- (axis cs:1e8,2.1)
  node[pos=0, above, font=\tiny, gray] {Chinchilla opt.};

\end{axis}
\end{tikzpicture}

\begin{tikzpicture}
\begin{axis}[
  width=0.82\textwidth,
  height=5.5cm,
  xmode=log, ymode=log,
  xlabel={Scale (compute / parameters / data)},
  ylabel={Loss},
  xlabel style={font=\small},
  ylabel style={font=\small},
  xmin=1e2, xmax=1e10,
  ymin=0.5, ymax=10,
  grid=major,
  grid style={gray!25, very thin},
  legend pos=north east,
  legend style={font=\small},
  tick label style={font=\tiny},
  title style={font=\small\bfseries},
]
% Compute scaling curve
\addplot[hitblue, very thick, domain=1e2:1e10, samples=50]
  {8 * x^(-0.05)};
\addlegendentry{Compute}

% Parameters scaling curve
\addplot[hitgold, very thick, dashed, domain=1e2:1e10, samples=50]
  {7 * x^(-0.045)};
\addlegendentry{Parameters}

% Data scaling curve
\addplot[hitalert, very thick, dotted, domain=1e2:1e10, samples=50]
  {9 * x^(-0.06)};
\addlegendentry{Dataset size}

% Marker: GPT-2 region
\node[font=\tiny, hitblue, rotate=30] at (axis cs:1e4, 3.8) {GPT-2};
% Marker: GPT-3 region
\node[font=\tiny, hitblue, rotate=30] at (axis cs:1e8, 1.8) {GPT-3};
% Marker: Chinchilla-optimal
\draw[->, gray, thin] (axis cs:3e7,2.5) -- (axis cs:1e8,2.1)
  node[pos=0, above, font=\tiny, gray] {Chinchilla opt.};

\end{axis}
\end{tikzpicture}

\begin{tikzpicture}
\begin{axis}[
  width=0.82\textwidth,
  height=5.5cm,
  xmode=log, ymode=log,
  xlabel={Scale (compute / parameters / data)},
  ylabel={Loss},
  xlabel style={font=\small},
  ylabel style={font=\small},
  xmin=1e2, xmax=1e10,
  ymin=0.5, ymax=10,
  grid=major,
  grid style={gray!25, very thin},
  legend pos=north east,
  legend style={font=\small},
  tick label style={font=\tiny},
  title style={font=\small\bfseries},
]
% Compute scaling curve
\addplot[hitblue, very thick, domain=1e2:1e10, samples=50]
  {8 * x^(-0.05)};
\addlegendentry{Compute}

% Parameters scaling curve
\addplot[hitgold, very thick, dashed, domain=1e2:1e10, samples=50]
  {7 * x^(-0.045)};
\addlegendentry{Parameters}

% Data scaling curve
\addplot[hitalert, very thick, dotted, domain=1e2:1e10, samples=50]
  {9 * x^(-0.06)};
\addlegendentry{Dataset size}

% Marker: GPT-2 region
\node[font=\tiny, hitblue, rotate=30] at (axis cs:1e4, 3.8) {GPT-2};
% Marker: GPT-3 region
\node[font=\tiny, hitblue, rotate=30] at (axis cs:1e8, 1.8) {GPT-3};
% Marker: Chinchilla-optimal
\draw[->, gray, thin] (axis cs:3e7,2.5) -- (axis cs:1e8,2.1)
  node[pos=0, above, font=\tiny, gray] {Chinchilla opt.};

\end{axis}
\end{tikzpicture}

  \end{columns}
  \note{The power law means you can predict performance before spending compute. This changed how labs plan training runs.}
\end{frame}

\begin{frame}{Chinchilla Scaling: Compute-Optimal Training}
  \textbf{Hoffmann et al.\ (2022)}~\cite{hoffmann2022chinchilla} revisited scaling laws:
  \vspace{0.4em}
  \begin{tcolorbox}[hitbox, title=Chinchilla Rule of Thumb]
    For compute-optimal training, \textbf{scale tokens and parameters equally}:
    \[
      N_\text{opt} \approx D_\text{opt} \propto C^{0.5}
    \]
    $\approx$ 20 training tokens per model parameter.
  \end{tcolorbox}
  \vspace{0.4em}
  \begin{itemize}
    \item GPT-3 (175B params, 300B tokens) was \emph{under-trained} by this criterion.
    \item Chinchilla (70B params, 1.4T tokens) outperformed GPT-3 at lower compute.
    \item Implication: many published LLMs are compute-suboptimal.
  \end{itemize}
  \note{Ask students: does this mean bigger is not always better? Yes — if data is cheap, train smaller models longer.}
\end{frame}

% ══════════════════════════════════════════════════════════════════════════════
\section{Emergent Capabilities}
% ══════════════════════════════════════════════════════════════════════════════

\sectionframe{Emergent Capabilities}{Abilities that appear suddenly with scale}

\begin{frame}{What Are Emergent Capabilities?}
  \begin{tcolorbox}[hitbox, title=Definition (Wei et al.\ 2022~\cite{wei2022emergent})]
    An ability is \textbf{emergent} if it is not present in smaller models but is present in larger models
    and cannot be predicted by extrapolating the performance of smaller models.
  \end{tcolorbox}
  \vspace{0.4em}
  \textbf{Examples of emergent abilities:}
  \begin{itemize}
    \item Few-shot in-context learning (appears $\sim$GPT-3 scale)
    \item Multi-step arithmetic reasoning
    \item Code generation, translation to/from formal languages
    \item Theory-of-mind reasoning (contested)
  \end{itemize}
  \note{Important caveat from Schaeffer et al. (2023): some emergence may be an artifact of discontinuous metrics. Discuss with students.}
\end{frame}

\begin{frame}{Emergent Capabilities: Caveats and Debates}
  \begin{columns}[T]
    \column{0.52\textwidth}
      \textbf{Arguments for genuine emergence:}
      \begin{itemize}
        \item Sharp phase transitions in benchmark accuracy
        \item Not predictable from smooth loss curves
      \end{itemize}
      \vspace{0.5em}
      \textbf{Counterarguments (Schaeffer et al.\ 2023):}
      \begin{itemize}
        \item Emergence disappears with \emph{continuous metrics}
        \item Artifact of threshold-based evaluation
        \item Gradual improvement in log-probability space
      \end{itemize}
    \column{0.44\textwidth}
      \placeholder{0.95\textwidth}{3.5cm}{BIG-Bench: accuracy jumps at GPT-3+ scale}
  \end{columns}
  \note{This is an active research debate. The takeaway is: be skeptical of claims that capability X ``appeared'' at scale Y.}
\end{frame}

% ══════════════════════════════════════════════════════════════════════════════
\section{The Foundation Model Paradigm}
% ══════════════════════════════════════════════════════════════════════════════

\sectionframe{The Foundation Model Paradigm}{One model, many tasks}

\begin{frame}{What Is a Foundation Model?}
  \begin{tcolorbox}[hitbox, title=Definition (Bommasani et al.\ 2021~\cite{bommasani2021foundation})]
    A \textbf{foundation model} is a model trained on broad data at scale that can be
    \emph{adapted} (via fine-tuning, prompting, or in-context learning) to a wide range of
    downstream tasks.
  \end{tcolorbox}
  \vspace{0.4em}
  \begin{columns}[T]
    \column{0.50\textwidth}
      \textbf{Key properties:}
      \begin{itemize}
        \item Trained on massive, diverse corpora
        \item General-purpose representations
        \item Adapt via prompting or fine-tuning
        \item Examples: GPT-4, Llama 2, Gemini, Claude
      \end{itemize}
    \column{0.45\textwidth}
      \placeholder{0.95\textwidth}{3.5cm}{Foundation model → many downstream tasks}
  \end{columns}
  \note{The key word is ``adaptation'' not retraining. This is what separates foundation models from earlier task-specific models.}
\end{frame}

\begin{frame}{Foundation Models vs.\ Traditional ML}
  \begin{center}
  \begin{tabular}{lll}
    \toprule
    & \textbf{Traditional ML} & \textbf{Foundation Model} \\
    \midrule
    Training data & Task-labeled & Broad, self-supervised \\
    Model size & Small–medium & Billions of parameters \\
    Adaptation & Retrain from scratch & Prompt / fine-tune \\
    Tasks & One model per task & One model, many tasks \\
    Compute & Low & Very high (pretraining) \\
    Interpretability & Often high & Often low \\
    \bottomrule
  \end{tabular}
  \end{center}
  \note{This table is a simplified view. In practice, both paradigms coexist, and foundation models are often fine-tuned.}
\end{frame}

% ══════════════════════════════════════════════════════════════════════════════
\section{Opportunities and Risks}
% ══════════════════════════════════════════════════════════════════════════════

\sectionframe{Opportunities and Risks}{The two sides of the foundation model coin}

\begin{frame}{Opportunities}
  \begin{columns}[T]
    \column{0.48\textwidth}
      \textbf{Scientific and economic value:}
      \begin{itemize}
        \item Democratize AI capabilities
        \item Accelerate scientific research (biology, chemistry, medicine)
        \item Reduce per-task development cost
        \item Multimodal understanding
      \end{itemize}
    \column{0.48\textwidth}
      \textbf{Practical applications:}
      \begin{itemize}
        \item Code generation and debugging
        \item Medical report generation
        \item Legal document analysis
        \item Educational tutoring
        \item Scientific hypothesis generation
      \end{itemize}
  \end{columns}
  \note{Bommasani et al. go into extensive detail. Students should read the original paper.}
\end{frame}

\begin{frame}{Risks and Challenges}
  \begin{tcolorbox}[alertbox, title=Key Risks (Bommasani et al.\ 2021)]
    \begin{itemize}
      \item \textbf{Bias amplification}: Models inherit and amplify biases in training data.
      \item \textbf{Misuse}: Disinformation, spam, phishing at scale.
      \item \textbf{Centralization}: Power concentrated in few organizations.
      \item \textbf{Homogenization}: Single model for all tasks → correlated failures.
      \item \textbf{Environmental cost}: Massive CO$_2$ footprint of pretraining.
      \item \textbf{Opacity}: Reasoning is often not interpretable.
    \end{itemize}
  \end{tcolorbox}
  \note{This sets the motivation for the safety and alignment lectures in Weeks 12 and 13.}
\end{frame}

% ══════════════════════════════════════════════════════════════════════════════
\section{Paper Discussion}
% ══════════════════════════════════════════════════════════════════════════════

\begin{frame}{Paper Discussion}
  \papercite{On the Opportunities and Risks of Foundation Models}
    {Bommasani et al.\ (Stanford HAI + 100 co-authors)}
    {arXiv 2021 — 200-page report}
  \vspace{0.5em}
  \textbf{Key contributions:}
  \begin{itemize}
    \item Coined the term \emph{foundation model}.
    \item Comprehensive taxonomy: capabilities, applications, societal considerations.
    \item Section on \emph{homogenization}: if foundation models fail, they fail everywhere.
  \end{itemize}
  \vspace{0.3em}
  \textbf{Discussion question:} Is the ``foundation'' metaphor helpful or misleading?\\
  What does it imply about the models' role in future AI systems?
  \note{Students should read the abstract and Section 1 (Introduction) and Section 4 (Societal Considerations) before this lecture.}
\end{frame}

% ══════════════════════════════════════════════════════════════════════════════
\section{Lab / Demo}
% ══════════════════════════════════════════════════════════════════════════════

\begin{frame}{Lab Idea: Probing Scaling Behavior}
  \begin{tcolorbox}[hitbox, title=Lab 1 — Scaling and Emergence]
    \textbf{Goal:} Observe how model capability changes with scale.\\[4pt]
    \textbf{Tools:} Hugging Face \texttt{transformers}, GPT-2 (117M, 345M, 774M, 1.5B)
    \begin{enumerate}
      \item Measure perplexity on WikiText-103 for each model size.
      \item Evaluate on a simple arithmetic task (3-digit addition).
      \item Plot accuracy vs.\ model size on a log-linear scale.
      \item Identify the scale at which the task ``emerges.''
    \end{enumerate}
    \textbf{Deliverable:} 1-page report + plots.
  \end{tcolorbox}
  \note{This lab can be done in 2–3 hours on a single GPU or Google Colab.}
\end{frame}

% ══════════════════════════════════════════════════════════════════════════════
\section{Discussion \& Summary}
% ══════════════════════════════════════════════════════════════════════════════

\begin{frame}{Discussion Questions}
  \begin{enumerate}
    \item The Chinchilla paper suggests many LLMs are \emph{under-trained} relative to their size.
          What are the practical implications for organizations that deploy these models?
    \item Emergent capabilities make it hard to \emph{predict} model behavior.
          How should this affect AI governance and deployment decisions?
    \item Is the ``foundation'' metaphor in ``foundation models'' accurate?
          What are the risks of building critical systems on top of a single base model?
    \item Should foundation model pretraining runs require environmental impact statements?
  \end{enumerate}
\end{frame}

\begin{frame}{Summary}
  \begin{itemize}
    \item \textbf{Transfer learning} enabled reuse of representations across tasks.
    \item \textbf{Scaling laws} (Kaplan 2020, Hoffmann 2022) show predictable power-law
          improvements and guide optimal compute allocation.
    \item \textbf{Emergent capabilities} appear at scale and are not fully predictable.
    \item \textbf{Foundation models} are trained broadly, adapted cheaply —
          a paradigm shift from task-specific ML.
    \item They bring both \textbf{opportunities} (democratization, science) and
          \textbf{risks} (bias, misuse, centralization).
  \end{itemize}
  \vspace{0.4em}
  \textbf{Next week:} Transformer architectures revisited — the machinery behind foundation models.
\end{frame}

\begin{frame}[allowframebreaks]{Suggested Reading}
  \nocite{bommasani2021foundation, kaplan2020scaling, hoffmann2022chinchilla,
          wei2022emergent, brown2020language, devlin2019bert}
  \bibliography{../../references}
\end{frame}

\end{document}
